\chapter{Methodology}\label{chap:Methodology}

This chapter describes the methodology used to design and implement Task~Flow.
It begins with a structured requirements analysis, then presents the high-level
system architecture, followed by detailed accounts of the data layer, backend,
frontend, and desktop-packaging implementations.
The chapter closes with a discussion of the testing strategy applied throughout
development.

% ─────────────────────────────────────────────────────────────────────────────
\section{Requirements Analysis}\label{sec:meth_requirements}
% ─────────────────────────────────────────────────────────────────────────────

\subsection*{Functional Requirements}

The following functional requirements were identified from the project objectives stated
in Section~\ref{sec:intro_aims}:

\begin{enumerate}
    \item \textbf{FR-01 User Authentication.}
          The system shall support user registration, login, and password reset via
          email.
          Passwords shall be stored as bcrypt hashes; sessions shall be managed with
          signed JWT tokens valid for 24~hours.

    \item \textbf{FR-02 Task Management.}
          Users shall be able to create, read, update, and delete tasks.
          Each task shall carry a title, optional description, priority (Low, Medium,
          High), status (To~Do, In~Progress, Review, Completed, Overdue), optional due
          date, star flag, and project assignment.

    \item \textbf{FR-03 Multiple Task Views.}
          Tasks shall be displayable in five views: a smart-grouped default list, a
          Kanban board, a tabular grid, a Gantt timeline, and a calendar time-grid.

    \item \textbf{FR-04 Project and Team Management.}
          Users shall be able to create projects, invite team members, and assign tasks
          to projects.
          Team invitations shall be transmitted via MongoDB to support cross-device
          reception.

    \item \textbf{FR-05 Messaging.}
          Users shall be able to exchange one-to-one direct messages and participate in
          group chats.
          Both conversation types shall support file attachments up to 50~MB.

    \item \textbf{FR-06 Notifications and Reminders.}
          The system shall deliver real-time in-application notifications over SignalR
          and optional email reminders at user-configured times.
          Automated warnings shall fire 24~hours and 1~hour before task due dates.

    \item \textbf{FR-07 Calendar Events.}
          Users shall be able to create, edit, and delete calendar events with start and
          end times, an optional colour, and an optional meeting link.

    \item \textbf{FR-08 AI Chatbot.}
          A conversational AI assistant shall be available in three modes: general
          chat, code generation, and document scanning (OCR).
          Responses shall be streamed token-by-token via SSE.
          Conversation history shall be persisted between sessions.

    \item \textbf{FR-09 Offline Operation.}
          All features in FR-01 through FR-07 shall function without network access.
          Locally generated changes shall be queued and replayed to MongoDB when
          connectivity is restored.

    \item \textbf{FR-10 Dashboard.}
          A dashboard shall display summary statistics (active tasks, in-progress count,
          project count, team member count), a task-trend indicator, tasks due this week,
          a recent-activity feed, and an embedded calendar widget.

    \item \textbf{FR-11 Settings.}
          Users shall be able to update their profile (name, avatar, company, country,
          phone, timezone), change their password, and manage notification preferences.
\end{enumerate}

\subsection*{Non-Functional Requirements}

\begin{enumerate}
    \item \textbf{NFR-01 Offline Availability.}
          The application shall start and remain fully operable on a machine with no
          network access.

    \item \textbf{NFR-02 Security.}
          All API endpoints shall require a valid JWT.
          Passwords shall not be stored in recoverable form.
          The application shall not expose SQL injection vectors.

    \item \textbf{NFR-03 Cross-Platform Distribution.}
          The packaged installer shall run on Windows (x64), macOS (x64/ARM), and Linux
          (x64) without requiring the user to install any additional runtime.

    \item \textbf{NFR-04 Responsiveness.}
          UI interactions shall produce visible feedback within 200~ms under normal
          conditions.

    \item \textbf{NFR-05 Extensibility.}
          The backend shall use a repository and service pattern such that new entity
          types can be added without modifying existing controller logic.
\end{enumerate}

% ─────────────────────────────────────────────────────────────────────────────
\section{System Architecture Overview}\label{sec:meth_arch}
% ─────────────────────────────────────────────────────────────────────────────

Task~Flow is structured as a three-tier desktop application embedded in an Electron shell.

\textbf{Tier 1 --- Presentation (React~+~TypeScript).}
The user interface is a single-page application (SPA) built with React~19 and
TypeScript~6.
It communicates with the backend exclusively via HTTP REST calls and a persistent
SignalR WebSocket connection.
Webpack~5 bundles all TypeScript and Tailwind CSS sources into a set of static assets
served from \texttt{wwwroot/} by the Kestrel HTTP server.

\textbf{Tier 2 --- Application Logic (ASP.NET Core 10).}
The backend exposes 14 REST API controllers grouped under \texttt{/api/}, one SignalR
hub (\texttt{/hubs/notifications}), and four background services.
It follows a layered architecture: controllers delegate to service classes, which in
turn use repository objects (EF Core \texttt{DbContext} queries) for persistence.
AutoMapper projects entity objects onto DTOs at the service boundary, preventing
internal model details from leaking into the API response.
A global \texttt{ExceptionHandlingMiddleware} catches unhandled exceptions and returns
a uniform JSON error envelope, ensuring no stack trace is ever exposed to the client.

\textbf{Tier 3 --- Data (SQLite + MongoDB).}
SQLite serves as the primary, always-available local store.
MongoDB is an optional cloud relay used for cross-device features.
The application starts and runs normally even when the MongoDB connection string is
empty or the server is unreachable.

\textbf{Electron Shell.}
The Electron main process (\texttt{public/electron.js}) launches the compiled .NET
backend as a child process.
It monitors the child's standard output stream and waits for the line
\texttt{TASKFLOW\_BACKEND\_READY:\{address\}} before opening the browser window and
pointing it at the dynamically assigned \texttt{http://127.0.0.1:\{port\}} address.
This design ensures the UI is never shown before the API is ready to serve requests.

\textbf{Request Lifecycle.}
A typical authenticated request follows the path:
React component $\to$ \texttt{fetch()} with \texttt{Authorization: Bearer \{token\}}
header $\to$ Kestrel $\to$ JWT middleware validates signature and extracts claims $\to$
controller action $\to$ service $\to$ EF Core $\to$ SQLite.
If the operation mutates a syncable entity, the service additionally calls
\texttt{IMirrorService.MirrorAsync()}, which either performs a direct MongoDB upsert
(online) or enqueues a \texttt{SyncOutboxEntry} (offline).
Real-time events are broadcast to the connected React client via the SignalR hub
independently of the HTTP response path.

% ─────────────────────────────────────────────────────────────────────────────
\section{Data Layer Design}\label{sec:meth_data}
% ─────────────────────────────────────────────────────────────────────────────

\subsection*{The ISyncableEntity Interface}

Entities that must be mirrored to MongoDB implement the \texttt{ISyncableEntity}
interface, which exposes three properties:

\begin{itemize}
    \item \texttt{SyncId} (GUID) --- a stable, globally unique identifier assigned at
          creation time.
          It is used as the MongoDB document \texttt{\_id}, enabling idempotent upserts:
          re-running the same outbox entry twice produces the same document, not a
          duplicate.
    \item \texttt{UpdatedAt} (DateTime) --- the wall-clock timestamp of the most recent
          change, used as the arbitration key in last-write-wins conflict resolution.
    \item \texttt{IsSynced} (bool) --- a dirty flag that indicates whether the local
          record has been propagated to MongoDB.
\end{itemize}

The \texttt{AppDbContext.SaveChangesAsync()} override automatically stamps
\texttt{UpdatedAt} to \texttt{DateTime.UtcNow} and assigns a new \texttt{SyncId} for
newly added entities, eliminating boilerplate from every service call.
It also sets \texttt{IsSynced = false} for every modified entity to ensure that
outstanding changes are detected by the sync engine.

\subsection*{SQLite Entity Model}

The application manages eighteen entity types in SQLite, grouped as follows:

\begin{itemize}
    \item \textbf{Identity}: \texttt{AppUser}
    \item \textbf{Work}: \texttt{TaskItem}, \texttt{TaskComment}, \texttt{Project},
          \texttt{ProjectMember}
    \item \textbf{Team \& Social}: \texttt{Team}, \texttt{TeamMember},
          \texttt{LocalTeamMember}, \texttt{LocalInvitation}
    \item \textbf{Messaging}: \texttt{Message}, \texttt{GroupChat},
          \texttt{GroupChatMember}, \texttt{GroupMessage}
    \item \textbf{Time}: \texttt{CalendarEvent}, \texttt{Reminder}
    \item \textbf{Alerts}: \texttt{Notification}
    \item \textbf{AI}: \texttt{ChatbotConversation}, \texttt{ChatbotMessage}
    \item \textbf{Sync}: \texttt{SyncOutboxEntry}
\end{itemize}

Unique indexes on \texttt{SyncId} for the four core syncable types (Project,
TaskItem, Notification, Reminder) prevent duplicate inserts if the application
is restarted mid-sync.
\texttt{LocalInvitation} and \texttt{LocalTeamMember} are read-only mirrors of MongoDB
documents pulled on login; their \texttt{MongoId} field serves as the local primary key.

\subsection*{MongoDB Collection Design}

Five collections are maintained in the \texttt{TaskFlow} MongoDB database:

\begin{table}[H]
\centering
\caption{MongoDB collections and their purpose.}
\label{tab:mongo_collections}
\begin{tabular}{lll}
\toprule
\textbf{Collection} & \textbf{Key Index} & \textbf{Purpose} \\
\midrule
\texttt{tasks}             & \texttt{\_id} (SyncId)          & Mirror of SQLite task rows \\
\texttt{notifications}     & \texttt{\_id} (SyncId)          & Mirror of SQLite notifications \\
\texttt{team\_invitations} & (recipientEmail, status)        & Cross-device team invites \\
\texttt{team\_members}     & (teamId, ownerEmail)            & Team membership roster \\
\texttt{user\_presence}    & email (unique)                  & Online/offline status \\
\bottomrule
\end{tabular}
\end{table}

% ─────────────────────────────────────────────────────────────────────────────
\section{Backend Implementation}\label{sec:meth_backend}
% ─────────────────────────────────────────────────────────────────────────────

\subsection{Authentication Service}\label{sec:meth_auth}

User registration follows a three-step process:
(1)~validate the input DTO using FluentValidation~\cite{aspnetcore2024};
(2)~hash the password with BCrypt.Net-Next using the default cost factor;
(3)~persist the \texttt{AppUser} record to SQLite and asynchronously mirror a safe
projection of the user to MongoDB, enabling other team members to look up the user by
email when sending invitations.

Login validates the submitted password against the stored hash using
\texttt{BCrypt.Verify()}.
On success, the service constructs a JWT containing ten claims
(\texttt{id}, \texttt{email}, \texttt{given\_name}, \texttt{family\_name},
\texttt{avatarUrl}, \texttt{company}, \texttt{country}, \texttt{phone},
\texttt{timezone}, and \texttt{jti}).
The token is signed with HMAC-SHA256 using the key resolved from the
\texttt{TASKFLOW\_JWT\_KEY} environment variable (falling back to the
\texttt{Jwt:Key} configuration value, or an auto-generated 32-byte random hex string
if neither is present).
A fire-and-forget call to \texttt{UserDataSyncService.PullForUserAsync()} pulls the
latest MongoDB data for that user immediately after login, refreshing local copies of
team memberships and pending invitations.

Password reset is implemented via a time-limited token: a random URL-safe token is
stored in \texttt{AppUser.ResetToken} alongside a \texttt{ResetTokenExpiry} timestamp
and emailed to the user using MailKit over Gmail SMTP (port~587, STARTTLS).
The reset endpoint validates the token, checks expiry, and updates the password hash.

\subsection{Task and Project Management}\label{sec:meth_tasks}

\texttt{TaskItem} is the central entity of the system.
Each task belongs to exactly one \texttt{AppUser} (its assignee) and optionally to
one \texttt{Project}.
The \texttt{Priority} enum defines three levels: \texttt{Low} (0), \texttt{Medium} (1),
and \texttt{High} (2).
The \texttt{TaskItemStatus} enum defines five states: \texttt{Todo},
\texttt{InProgress}, \texttt{Review}, \texttt{Completed}, and \texttt{Overdue}.
The \texttt{Overdue} state is not set manually; it is assigned automatically by the
\texttt{DueDateWarningService} background service when the due date passes.

All \texttt{TaskService} operations validate ownership by comparing the
\texttt{AssigneeId} of the fetched entity against the \texttt{userId} extracted from
the JWT claim \texttt{ClaimTypes.NameIdentifier}, returning a 403 response for
unauthorised access attempts.
Text search is executed at the database layer using EF Core \texttt{.Where()} with
\texttt{string.Contains()} predicates (compiled to SQL \texttt{LIKE} clauses),
avoiding in-memory filtering of large datasets.
AutoMapper projects \texttt{TaskItem} entities onto \texttt{TaskItemDto} response
objects, exposing only the fields needed by the frontend.

After any create or update operation, the service invokes
\texttt{IMirrorService.MirrorAsync(entity, operationName)}.
\texttt{MirrorService} checks \texttt{ConnectivityService.IsEffectivelyOnline}: if
online, it calls \texttt{MongoService.UpsertDocumentBySyncIdAsync()} directly; if
offline, it serialises the entity payload to JSON and inserts a
\texttt{SyncOutboxEntry} into the local database with status \texttt{Pending}.

\subsection{Real-Time Notification System}\label{sec:meth_notifications}

\texttt{NotificationHub} extends ASP.NET Core SignalR's \texttt{Hub} class.
On connection, each client is added to a user-specific SignalR group named
\texttt{user\_\{userId\}}.
The hub reads the JWT from the \texttt{access\_token} query string parameter (required
because SignalR's JavaScript client cannot set HTTP headers during the WebSocket
upgrade).
The unread notification count is pushed to the client immediately upon connection.

The system defines fifteen notification types (e.g., \texttt{TaskAssigned},
\texttt{TaskCompleted}, \texttt{MessageReceived}, \texttt{TeamInvitation},
\texttt{TaskOverdue}) and four priority levels (\texttt{Low}, \texttt{Normal},
\texttt{High}, \texttt{Urgent}).
Each notification includes an \texttt{ActionUrl} field that the frontend uses to
navigate the user to the relevant resource on click.

Two background services supplement the notification hub:
\texttt{DueDateWarningService} runs on a 1-minute timer, queries all tasks with
approaching due dates, and fires both in-app notifications and (if configured) email
alerts.
An in-memory \texttt{\_sentNotifications} dictionary deduplicates warnings so that
the same task does not generate repeated notifications within the same application
session.
\texttt{ReminderProcessorService} runs on a 1-minute timer and fires user-configured
\texttt{Reminder} records whose \texttt{FireAt} time has passed and whose
\texttt{HasFired} flag is still false.

\subsection{Messaging System}\label{sec:meth_messages}

The messaging module provides two conversation types.

\textbf{Direct Messages.}
The \texttt{Message} entity stores a sender, a receiver, a text body, a read flag, and
separate soft-delete flags for each party (\texttt{IsDeletedBySender},
\texttt{IsDeletedByReceiver}).
This design ensures that a message deleted by one party remains visible to the other
until they also delete it.
File attachments (URL, original file name, MIME type, size in bytes) are stored as
nullable fields on the \texttt{Message} entity and are limited to 50~MB by Kestrel's
request body size limit.
The \texttt{MessagesController} sends a \texttt{MessageReceived} SignalR event to the
recipient's user group after persisting each message, enabling real-time delivery without
polling.

\textbf{Group Chats.}
\texttt{GroupChat} stores a name, description, and owner.
\texttt{GroupChatMember} is the join entity between a group and its participants.
\texttt{GroupMessage} mirrors the attachment model of direct messages and additionally
records whether the message has been edited (\texttt{IsEdited}).
The \texttt{GroupChatsController} broadcasts a \texttt{GroupMessageReceived} event
to a SignalR group named after the group chat ID.

\subsection{Offline Synchronisation Engine}\label{sec:meth_sync}

The synchronisation subsystem is composed of three cooperating components.

\textbf{ConnectivityService} is a singleton that maintains the effective connectivity
state of the application.
It exposes \texttt{IsOnline} (result of the last MongoDB ping), \texttt{IsManualOffline}
(user-toggled flag from the UI), and the derived \texttt{IsEffectivelyOnline} property
(\texttt{IsOnline \&\& !IsManualOffline}).
It also exposes \texttt{PendingSyncCount} (count of \texttt{SyncOutboxEntry} rows with
status \texttt{Pending} or \texttt{Failed}) and a \texttt{ConnectivityChanged} event
that subscribers can observe.

\textbf{OfflineSyncService} is the background worker that replays the outbox.
It runs a 10-second ping loop: on each tick it calls \texttt{MongoService.PingAsync()}
and updates \texttt{ConnectivityService.IsOnline}.
It also subscribes to \texttt{ConnectivityChanged}: when the state transitions from
offline to online, it waits 1500~ms (allowing the MongoDB connection to stabilise) and
then calls \texttt{ReplayOutboxAsync()}.
During replay, each \texttt{Pending} or \texttt{Failed} outbox entry is advanced to
\texttt{Processing}, the stored operation is executed against MongoDB, and the entry is
marked \texttt{Synced} on success or \texttt{Failed} (with an incremented attempt
counter and error message) on failure.
SignalR events (\texttt{SyncStarted}, \texttt{SyncProgress}, \texttt{SyncComplete}) are
broadcast to the client throughout the replay process, feeding the
\texttt{ConnectivityBar}'s animated progress indicator.

\textbf{BulkSyncStartupService} runs once at application startup.
It waits up to 90~seconds for MongoDB to become reachable, then performs a bulk upsert
of all SQLite entities implementing \texttt{ISyncableEntity}.
This ensures that a fresh installation, or a machine that was offline during the
previous session, fully reconciles its local state with MongoDB before accepting new
writes.

\subsection{AI Integration}\label{sec:meth_ai}

\texttt{MistralChatService} encapsulates all communication with the Mistral REST API.
It is constructed as a typed \texttt{HttpClient} singleton registered in the DI
container.
Three operational modes are supported, each targeting a different model and system
prompt:

\begin{itemize}
    \item \textbf{General mode} (\texttt{mistral-small-latest}) --- a general-purpose
          assistant that can answer questions, help with planning, and discuss any topic.
    \item \textbf{Coder mode} (\texttt{codestral-latest}) --- a code-specialised
          assistant with a system prompt that instructs the model to prefer code
          examples and to explain algorithmic choices.
    \item \textbf{Scanner mode} (\texttt{mistral-ocr-latest}) --- a vision-language
          mode that accepts base64-encoded images or PDF documents and returns extracted,
          structured text.
          PDFs are submitted to the \texttt{/v1/ocr} endpoint; images are embedded
          directly in the chat message payload.
\end{itemize}

\textbf{Streaming.}
The \texttt{ChatAsync()} method makes an HTTP POST request with
\texttt{stream: true} in the request body and reads the response using
\texttt{HttpCompletionOption.ResponseHeadersRead}, then processes the response body as
an \texttt{IAsyncEnumerable<string>} of SSE lines.
Each line prefixed with \texttt{data:} is deserialised as a chat-completion delta, and
the token content is written immediately to the ASP.NET Core SSE response stream using
\texttt{Response.WriteAsync()}.
This pipeline adds negligible latency between the Mistral API producing a token and the
React frontend rendering it.

\textbf{History window.}
The service loads the last 20 messages from the \texttt{ChatbotMessage} table before
constructing the API request, providing the model with sufficient context while bounding
the token cost of each request.
After streaming completes, the assistant's full response is assembled from all received
deltas and persisted as a new \texttt{ChatbotMessage} row with \texttt{Role = "assistant"}.

\textbf{Title generation.}
When a new conversation is started, a follow-up call to \texttt{GenerateTitleAsync()}
asks the model to produce a short, descriptive title from the first user message.
This title is stored in \texttt{ChatbotConversation.Title} and displayed in the
conversation list sidebar.

% ─────────────────────────────────────────────────────────────────────────────
\section{Frontend Implementation}\label{sec:meth_frontend}
% ─────────────────────────────────────────────────────────────────────────────

\subsection{Application Structure}\label{sec:meth_fe_structure}

The React application is organised as a flat module tree rooted at
\texttt{ReactApp/}.
React Router~7 manages client-side navigation with a route table of seventeen pages.
All routes except \texttt{/login}, \texttt{/signup}, \texttt{/forgot-password}, and
\texttt{/reset-password} are wrapped in an \texttt{<AuthGuard>} component that
redirects unauthenticated users to the login page.

Styling is managed entirely with Tailwind CSS~4 utility classes; no custom CSS files
exist in the project.
The Lucide React icon library provides consistent SVG icons throughout the UI.
Radix UI primitives (\texttt{@radix-ui/react-dropdown-menu}) are used for accessible
dropdown menus.
The \texttt{react-markdown} library with the \texttt{remark-gfm} plugin renders
assistant responses as formatted Markdown, supporting code blocks, tables, and inline
formatting.

Webpack~5 transpiles TypeScript via \texttt{ts-loader} and PostCSS processes Tailwind
classes; the output bundles are written to \texttt{wwwroot/dist/} and served directly
by Kestrel.

\subsection{Task Visualisation Views}\label{sec:meth_views}

All five views consume the same \texttt{TaskItem[]} array from a shared context and
differ only in their rendering and interaction model.

\textbf{Default View.}
Tasks are partitioned into five temporal groups computed client-side: Overdue (past due
date), Today (due today), This Week (due within 7~days), Later (due after 7~days), and
Completed.
Each group is rendered as a collapsible section with a task count badge.
This grouping replicates the mental model of a physical in-tray sorted by urgency.

\textbf{Kanban View.}
Five columns correspond to the five task statuses.
Each column header is coloured to convey urgency (red for Overdue, amber for In
Progress, etc.) and displays the column's task count.
Tasks are rendered as drag-and-drop cards using the HTML5 Drag and Drop API; dropping a
card in a different column fires a \texttt{PATCH /api/tasks/\{id\}} request to update
the task's status.

\textbf{Table View.}
Tasks are displayed in a columnar grid with columns for title, status, priority, due
date, project, and actions.
Status and priority are rendered as colour-coded badge components.
Each row has an expand toggle that reveals a sub-task scaffold row (UI-level only;
sub-task backend persistence is a planned future feature).

\textbf{Gantt View.}
A full-calendar-year timeline is rendered as a horizontal structure.
Each task is represented as a bar positioned and sized according to its due date.
A priority colour legend (red = High, amber = Medium, green = Low) allows at-a-glance
workload assessment.
The current date is highlighted with a vertical marker.

\textbf{Calendar View.}
A weekly time-grid spanning 7:00~AM to 8:00~PM displays tasks and calendar events as
positioned blocks.
A dark mini-calendar sidebar allows week navigation.
Tasks without a specific time are shown as all-day entries at the top of the grid.

\subsection{Connectivity Bar and Offline UX}\label{sec:meth_connbar}

\texttt{ConnectivityBar.tsx} is a fixed-position banner rendered at the top of the
application layout.
It subscribes to real-time SignalR events (\texttt{ConnectivityChanged},
\texttt{SyncStarted}, \texttt{SyncProgress}, \texttt{SyncComplete}) via the
\texttt{useConnectivity()} custom hook.

When the application is offline, the banner shows in amber with an offline icon and the
count of pending sync operations.
When a sync pass is in progress, the banner shows in blue and renders an animated
progress bar that advances as each outbox entry is processed.
When connectivity is healthy and no sync is in progress, the banner is hidden entirely
to avoid distracting the user.
A \texttt{toggleManualOffline} button allows developers and testers to simulate offline
mode without physically disconnecting the network.

\subsection{AI Chatbot Interface}\label{sec:meth_chatbot}

\texttt{Chatbot.tsx} is the most complex single component in the application.
Three toggle buttons at the top of the chat panel select the active mode (General,
Coder, Scanner); the active mode determines which API model is used on the backend.

When no conversation is active, a grid of starter-prompt buttons provides suggested
questions appropriate for each mode, helping new users understand the assistant's
capabilities.

Messages are submitted via a form that supports both text and a single file attachment
(triggered by a paperclip icon).
The component opens an \texttt{EventSource} connection to the SSE endpoint for each
new message and appends incoming token chunks to the current assistant message bubble
as they arrive.
When the SSE stream closes, the final assembled message is rendered through
\texttt{react-markdown} with syntax-highlighted code blocks.

A collapsible sidebar lists all saved conversations, sorted by most-recently-updated.
Each conversation displays its AI-generated title and a relative timestamp.
Clicking a conversation loads its full message history from the backend.

% ─────────────────────────────────────────────────────────────────────────────
\section{Desktop Packaging}\label{sec:meth_electron}
% ─────────────────────────────────────────────────────────────────────────────

The production build of Task~Flow is a self-contained desktop installer that bundles
three components: the Electron shell, the compiled .NET backend, and the Webpack-built
React frontend assets.

\textbf{Electron main process.}
\texttt{public/electron.js} is the entry point for the desktop application.
On startup it spawns the compiled .NET backend executable as a child process, capturing
its standard output.
The main process watches for the line matching \texttt{TASKFLOW\_BACKEND\_READY:\{address\}}
to appear.
Once detected, it extracts the port number, constructs the backend URL, and calls
\texttt{mainWindow.loadURL()} to load the React SPA from the local Kestrel server.
If the backend process exits unexpectedly, the main process terminates the Electron
application with an error dialog.

\textbf{CORS configuration.}
Because Electron loads the frontend from a \texttt{file://} origin (for the production
build) or from \texttt{localhost} (during development), the ASP.NET Core CORS policy
uses \texttt{SetIsOriginAllowed(\_ => true)} combined with \texttt{AllowCredentials()}
to permit the SignalR connection from either origin without hard-coding specific URLs.

\textbf{JWT key management.}
The Kestrel server listens on \texttt{http://127.0.0.1:0} (a random available port).
The JWT signing key is resolved in order of priority:
(1)~\texttt{TASKFLOW\_JWT\_KEY} environment variable;
(2)~\texttt{Jwt:Key} value in \texttt{appsettings.json};
(3)~an automatically generated 32-byte random hex string persisted for the duration of
the process.
This means the application works out-of-the-box without any configuration while still
allowing a stable key to be injected for production deployments.

\textbf{electron-builder configuration.}
\texttt{electron-builder.json} configures the following target artefacts:

\begin{itemize}
    \item \textbf{Windows} --- NSIS installer with \texttt{oneClick: false}, allowing
          the user to choose the installation directory and optionally create desktop
          and Start Menu shortcuts.
    \item \textbf{macOS} --- DMG image and ZIP archive (for notarisation-free
          distribution).
    \item \textbf{Linux} --- AppImage (portable, no installation required) and DEB
          package (for Debian/Ubuntu-based distributions).
\end{itemize}

The application ID is \texttt{com.taskflow.desktop}.
The .NET backend self-contained publish output (\texttt{win-x64}, \texttt{osx-x64},
\texttt{linux-x64}) is included as an extra resource so that electron-builder packages
the correct native binary for each target platform.

% ─────────────────────────────────────────────────────────────────────────────
\section{Testing Strategy}\label{sec:meth_testing}
% ─────────────────────────────────────────────────────────────────────────────

Testing was conducted at three levels throughout the development process.

\textbf{Unit-level verification.}
Service-layer methods were verified by constructing the service with an in-memory
SQLite database (using EF Core's \texttt{UseInMemoryDatabase} provider) and asserting
the expected state after each operation.
JWT claim extraction and bcrypt hash/verify round-trips were tested in isolation.

\textbf{API integration testing.}
Each of the 14 controllers was exercised against a running Kestrel instance using
\texttt{HttpClient} requests constructed with valid and invalid JWT tokens.
The tests verified: (a)~that protected endpoints return~401 when no token is supplied;
(b)~that ownership checks return~403 for resources belonging to a different user;
and (c)~that create, read, update, and delete operations produce the expected HTTP
status codes and response bodies.

\textbf{Offline synchronisation testing.}
The offline sync pipeline was validated using the \texttt{ConnectivityBar}'s manual
offline toggle.
With \texttt{IsManualOffline = true}, a series of task create and update operations was
performed; the outbox was inspected to confirm that \texttt{SyncOutboxEntry} rows were
created with status \texttt{Pending}.
On toggling back online, the sync replay was observed via the progress bar and the
MongoDB collection was inspected to confirm all entries were promoted to
\texttt{Synced}.

\textbf{Cross-platform packaging.}
The NSIS installer was tested on Windows~11 and Windows~10.
The macOS DMG was verified on macOS~14 (Sonoma) with both Intel and Apple Silicon
(via Rosetta~2 emulation).
The AppImage was tested on Ubuntu~22.04.
In all cases the application launched, authenticated, and performed offline task
creation and synchronisation without errors.
